\documentclass[a4paper]{article}
\usepackage{graphicx}
\usepackage{amsmath,amssymb}
\usepackage{hyperref}
\usepackage{minted}
\usepackage{multirow}
\usepackage{float}
\usepackage{todonotes}
\usepackage{forest}
\usepackage{xstring}
\input{/home/donkarlo/Dropbox/repo/nd_language_ai_project/src/nd_language_ai/formal/computer/programming/tex/latex/code_clips/external_dot.tex}
\title{Ermittlung des Teilergebnisses „Schreiben“}
\date{}

\begin{document}
    \maketitle
    
    \textbf{Deutsch:} Entscheidend für den Prüfungserfolg sind die Teilergebnisse folgender Subtests.
    \textbf{English:} The partial results of the following subtests are decisive for success in the exam.
    
    \textbf{Deutsch:} Die Bewertung des Subtests „Schreiben“ erfolgt durch lizenzierte Bewerter und Bewerterinnen gemäß den Bewertungskriterien. Bei den Stufen B1 und A2 besteht die Möglichkeit zu unterscheiden, ob die Kriterien „gut erfüllt“ oder „erfüllt“ wurden. Wichtig ist aber stets, dass die Bewerter und Bewerterinnen ihr Urteil kriterienbasiert und nicht nach Punktwerten fällen.
    
    \textbf{English:} The evaluation of the subtest “Writing” is carried out by licensed examiners according to the evaluation criteria. At levels B1 and A2 it is possible to distinguish whether the criteria are “well fulfilled” or “fulfilled”. However, it is always important that examiners make their judgement based on criteria rather than purely on point values.
    
    \begin{table}[H]
        \centering
        \begin{tabular}{|l|c|c|c|c|c|c|}
            \hline
            & \multicolumn{2}{c|}{B1} & \multicolumn{2}{c|}{A2} & A1 & 0 \\
            \cline{2-7}
            & gut erfüllt      & erfüllt     & gut erfüllt      & erfüllt     & erfüllt     &     \\
            & (well fulfilled) & (fulfilled) & (well fulfilled) & (fulfilled) & (fulfilled) &     \\
            \hline
            I Inhalt (Content)                                 & 5                & 4           & 3                & 2           & 1           & 0   \\
            II Kommunikative Gestaltung (Communicative design) & 5                & 4           & 3                & 2           & 1           & 0   \\
            III Korrektheit (Correctness)                      & 5                & 4           & 3                & 2           & 1           & 0   \\
            IV Wortschatz (Vocabulary)                         & 5                & 4           & 3                & 2           & 1           & 0   \\
            \hline
            Summe (Total)                                      & 20               & 16          & 12               & 8           & 4           & 0   \\
            Anteil (Percentage)                                & 100\%            & 80\%        & 60\%             & 40\%        & 20\%        & 0\% \\
            \hline
        \end{tabular}
    \end{table}
    
    \textbf{Deutsch:} Für das Erreichen der Stufen A2 und B1 gilt:
    \textbf{English:} The following score ranges apply for reaching levels A2 and B1:
    
    \begin{table}[H]
        \centering
        \begin{tabular}{|c|c|}
            \hline
            Punkte (Points) & Stufen nach GER (CEFR Level) \\
            \hline
            15--20          & B1                           \\
            7--14           & A2                           \\
            0--6            & unter A2 (below A2)          \\
            \hline
        \end{tabular}
    \end{table}
    % \bibliographystyle{plain}
    % \bibliography{/home/donkarlo/Dropbox/repo/shared_research/bibliography/bibliography.bib}
\end{document}